% ============================================================
\setpartpreamble{\partintro
  {Testing, versioning, packaging, and the checklist from A to Z.}
  {
    \item Chapter~\ref{ch:test-strategy}: Test strategy
    \item Chapter~\ref{ch:versioning-and-compatibility}: Versioning and compatibility
    \item Chapter~\ref{ch:packaging-and-documentation}: Packaging and documentation
    \item Chapter~\ref{ch:the-checklist-from-a-to}: The checklist from A to Z
  }}
\part{Quality and delivery}\label{part:VII}

% ------------------------------------------------------------
\chapter{Test strategy}\label{ch:test-strategy}

Plugins can be tested to a surprising extent without a database. Model building, type
mappings, SQL translation, migration SQL and snapshots only need a connection string, no
connection. The sample's test strategy has four levels:

\begin{datatable}{H{0.78}H{1.70}H{0.52}}
  \textbf{Level} & \textbf{Checks} & \textbf{Database} \\
  \tablehead
  Model tests & Annotations, conventions, validation, options, caching & no \\
  \tablerule
  SQL tests & Translations via \code{To\allowbreak{}Query\allowbreak{}String}, migration SQL
  via differ and generator & no \\
  \tablerule
  Snapshot tests & Code generator as a string and as a round trip (Roslyn) & no \\
  \tablerule
  Integration tests & Real behavior: RLS, roles, interceptors (Testcontainers) & yes
  (Docker) \\
\end{datatable}

\section{Test helpers}

Tests build many contexts with different \code{On\allowbreak{}Model\allowbreak{}Creating}.
Because EF caches the model per internal provider, every test must get its own provider;
that is what
\code{Enable\allowbreak{}Service\allowbreak{}Provider\allowbreak{}Caching(\allowbreak{}false)}
is for:

\begin{sourcecode}[tests/\allowbreak{}.\allowbreak{}.\allowbreak{}.\allowbreak{}/\allowbreak{}Test\allowbreak{}Model.\allowbreak{}cs (\allowbreak{}shortened)]{csharp}
public static TestContext Create(
    Action<TenancyOptionsBuilder>? tenancy = null,
    Action<ModelBuilder>? model = null,
    bool cacheServiceProvider = false,
    Action<DbContextOptionsBuilder>? configure = null)
{
    var builder = new DbContextOptionsBuilder<TestContext>()
        .UseNpgsql("Host=localhost;Database=unused")       // no connection needed
        .UseAcmeTenancy(tenancy)
        .EnableServiceProviderCaching(cacheServiceProvider);
    configure?.Invoke(builder);
    return new TestContext(builder.Options, model);
}

public static IReadOnlyList<string> MigrationSql(DbContext? source, DbContext target)
{
    var operations = target.GetService<IMigrationsModelDiffer>().GetDifferences(
        source?.GetService<IDesignTimeModel>().Model.GetRelationalModel(),
        target.GetService<IDesignTimeModel>().Model.GetRelationalModel());

    return target.GetService<IMigrationsSqlGenerator>()
        .Generate(operations, target.GetService<IDesignTimeModel>().Model)
        .Select(c => c.CommandText.Trim()).ToList();
}
\end{sourcecode}

\begin{important}[Use the design-time model]
  Always use the \textbf{design-time model} for migration tests
  (\code{IDesign\allowbreak{}Time\allowbreak{}Model}). \code{context.\allowbreak{}Model} is
  the lean runtime model; annotations that the annotation provider only passes on at design
  time are missing from it (failure case in
  chapter~\ref{ch:declarative-sql-from-annotations}).
\end{important}

\begin{note}[Why not SQLite or InMemory?]
  Tests against SQLite or the InMemory provider are tempting because they need no Docker.
  For a plugin they prove little: InMemory is not relational, so type mappings, translators
  and migrations never run; SQLite has a different SQL dialect, other type mappings and no
  Row Level Security. A provider-specific plugin is tested with its own provider: without a
  database (connection string only) for everything up to the SQL text, and against the real
  database for behavior.
\end{note}

\section{What is tested}

\begin{datatable}{H{0.75}H{1.25}}
  \textbf{Test} & \textbf{Ensures} \\
  \tablehead
  Caching with different options & Two contexts with different settings get different models
  (hash mistake from chapter~\ref{ch:the-options-extension}). \\
  \tablerule
  Invalid setting names & Validation including an injection attempt. \\
  \tablerule
  Type mapping without configuration & \code{Tenant\allowbreak{}Id} becomes \code{uuid};
  literal format; a different column type is not claimed. \\
  \tablerule
  \code{Current\allowbreak{}Tenant(\allowbreak{})} via
  \code{To\allowbreak{}Query\allowbreak{}String} & Translation, and that nothing is
  evaluated client-side (chapter~\ref{ch:preventing-client-side-evaluation}). \\
  \tablerule
  Explicit beats convention & An own policy with the same name is kept; the Fluent API
  overrides the attribute. \\
  \tablerule
  Warning as error & The opt-out warning becomes an error via
  \code{Configure\allowbreak{}Warnings}. \\
  \tablerule
  Validation & RLS on a derived type and invalid policies throw while building the model. \\
  \tablerule
  No phantom migration & An unchanged model produces not a single operation. \\
  \tablerule
  Migration SQL & Initial, change (DROP + CREATE), opt-out, custom operations. \\
  \tablerule
  SaveChanges logic & Setting, foreign tenant, reassignment, no tenant. \\
  \tablerule
  Snapshot & Fluent API instead of annotations; round trip without differences. \\
  \tablerule
  RLS in PostgreSQL & Isolation, no rows without a tenant, WITH CHECK applies. \\
  \tablerule
  Edge cases
  (chapters~\ref{ch:bulk-operations-and-the-limits}--\ref{ch:advanced-migrations}) &
  ExecuteUpdate and raw SQL against RLS, value generation, JSON columns, compiled model,
  idempotent script. \\
  \tablerule
  Composability & Decorator around the migrations generator and the composition matrix of
  two plugins (chapter~\ref{ch:registering-services}), EFCore.NamingConventions in both
  orders (chapter~\ref{ch:conventions}). \\
  \tablerule
  Annotation format & Serialized policies stay identical; tolerant reading
  (chapter~\ref{ch:versioning-and-compatibility}). \\
  \tablerule
  Context lifetimes & DbContext factory, background jobs, open connections
  (chapters~\ref{ch:interceptors}, \ref{ch:using-the-plugin-in-asp}). \\
  \tablerule
  Cross-tenant access & Reporting role reads all tenants but cannot write; the app role
  cannot switch RLS off (chapter~\ref{ch:row-level-security-in-production}). \\
\end{datatable}

\section{The snapshot round trip}

The most valuable test for design-time code: generate the snapshot, compile it with Roslyn,
build the model from it and diff it against the original. Any discrepancy between code
generator and Fluent API would otherwise produce a phantom migration for every user:

\begin{sourcecode}[tests/\allowbreak{}.\allowbreak{}.\allowbreak{}.\allowbreak{}/\allowbreak{}Snapshot\allowbreak{}Tests.\allowbreak{}cs (\allowbreak{}excerpt)]{csharp}
[Fact]
public void Snapshot_roundtrip_produces_no_differences()
{
    using var ctx = CreateContext();
    var snapshot = Compile(GenerateSnapshot(ctx));

    // As in EF's migrator: finalize and initialize
    var model = snapshot.Model;
    if (model is IMutableModel mutable) model = mutable.FinalizeModel();
    model = ctx.GetService<IModelRuntimeInitializer>()
        .Initialize(model, designTime: true, validationLogger: null);

    var differences = ctx.GetService<IMigrationsModelDiffer>().GetDifferences(
        model.GetRelationalModel(), ctx.DesignModel().GetRelationalModel());

    Assert.Empty(differences);
}
\end{sourcecode}

\section{Integration tests with Testcontainers}

The real security behavior can only be verified against PostgreSQL. The fixture starts a
container, creates the schema with \code{Ensure\allowbreak{}Created} (including RLS,
chapter~\ref{ch:how-migrations-work}), executes the custom migration operations for real and
connects the tests with a restricted role:

\begin{sourcecode}[tests/\allowbreak{}.\allowbreak{}.\allowbreak{}.\allowbreak{}Integration\allowbreak{}Tests/\allowbreak{}Postgres\allowbreak{}Fixture.\allowbreak{}cs (\allowbreak{}excerpt)]{csharp}
await owner.Database.EnsureCreatedAsync();          // RLS, policies, index

var commands = owner.GetService<IMigrationsSqlGenerator>().Generate(
[
    new CreateApplicationRoleOperation { Name = AppRole, CanLogin = true },
    new GrantTablePrivilegesOperation { Role = AppRole },
]);
await owner.GetService<IMigrationCommandExecutor>().ExecuteNonQueryAsync(
    commands, owner.GetService<IRelationalConnection>(),
    new MigrationExecutionState(), commitTransaction: true);
\end{sourcecode}

\begin{sourcecode}[tests/\allowbreak{}.\allowbreak{}.\allowbreak{}.\allowbreak{}Integration\allowbreak{}Tests/\allowbreak{}Row\allowbreak{}Level\allowbreak{}Security\allowbreak{}Tests.\allowbreak{}cs (\allowbreak{}excerpt)]{csharp}
[Fact]
public async Task Database_rejects_rows_of_other_tenant()
{
    // Bypass the interceptor: direct SQL. The policy's WITH CHECK must apply.
    var tenants = new MutableTenantProvider { CurrentTenantId = TenantId.New() };
    await using var db = _fixture.CreateAppContext(tenants);

    var ex = await Assert.ThrowsAsync<PostgresException>(() => db.Database.ExecuteSqlAsync(
        $"INSERT INTO \"Orders\" (\"TenantId\", \"Total\") VALUES ({Guid.NewGuid()}, 1)"));

    Assert.Equal(PostgresErrorCodes.InsufficientPrivilege, ex.SqlState);
}
\end{sourcecode}

\begin{tip}
  For providers and very deep plugins, EF Core provides the \textbf{specification tests}
  (\code{Microsoft.\allowbreak{}Entity\allowbreak{}Framework\allowbreak{}Core.\allowbreak{}Relational.\allowbreak{}Specification.\allowbreak{}Tests}).
  This is the same test suite EF and the providers themselves are tested against. The
  NetTopologySuite plugin in the EF repository shows how a plugin uses it.
\end{tip}

% ------------------------------------------------------------
\chapter{Versioning and compatibility}\label{ch:versioning-and-compatibility}

\section{Align versions with EF Core}

The plugin's major version follows EF Core's: version 10.x for EF Core 10. Users immediately
see which combination fits, and the plugin can adopt breaking changes with every EF major
version. Dependencies are limited to this major version:

\begin{sourcecode}[Directory.\allowbreak{}Packages.\allowbreak{}props (\allowbreak{}excerpt)]{xml}
<PackageVersion Include="Microsoft.EntityFrameworkCore" Version="[10.0.0,11.0.0)" />
<PackageVersion Include="Npgsql.EntityFrameworkCore.PostgreSQL" Version="[10.0.0,11.0.0)" />
\end{sourcecode}

\section{Dealing with internal APIs}

Plugins rarely get by without internal APIs: the sample derives from
\code{Npgsql\allowbreak{}Migrations\allowbreak{}Sql\allowbreak{}Generator},
\code{Npgsql\allowbreak{}Annotation\allowbreak{}Provider} and
\code{Npgsql\allowbreak{}Annotation\allowbreak{}Code\allowbreak{}Generator} and uses
\code{Npgsql\allowbreak{}Options\allowbreak{}Extension}. Such types can change even in minor
versions. EF reports their use with warning \textbf{EF1001}.

\begin{datatable}{H{0.80}H{1.20}}
  \textbf{Measure} & \textbf{Effect} \\
  \tablehead
  Look for a public alternative & Every internal API avoided is one risk less. \\
  \tablerule
  Bundle internal access & A few classes encapsulate it; a break affects one place. \\
  \tablerule
  Set
  \code{<\allowbreak{}No\allowbreak{}Warn>EF1001<\allowbreak{}/\allowbreak{}No\allowbreak{}Warn>}
  deliberately & Project-wide and documented, not scattered per line. \\
  \tablerule
  Limit the version range & Prevents users from getting untested combinations. \\
  \tablerule
  CI against previews & A build against the next EF preview shows breaks months earlier. \\
  \tablerule
  Round trip and SQL tests & Detect behavior changes the compiler does not see. \\
\end{datatable}

\begin{note}[Multi-targeting]
  Supporting several EF major versions at once is possible through multi-targeting, but
  doubles test effort and conditional compilation. For most plugins, one release branch per
  EF major version is simpler.
\end{note}

\section{Evolving annotation formats}

Everything a plugin writes into annotations ends up in its users' model snapshots, often
long before the next upgrade. If a new plugin version serializes the same configuration
differently (an additional JSON property, a different order, another enum spelling), the
differ compares two different strings. Every user then gets a migration that drops and
recreates every policy without any actual change. The serialized format is therefore part of
the public API, just like the methods.

\begin{datatable}{H{0.79}H{1.21}}
  \textbf{Rule} & \textbf{How} \\
  \tablehead
  Pin the format with a test & A golden-value test compares the serialized form character by
  character. \\
  \tablerule
  Read tolerantly & Missing properties become defaults, unknown ones are ignored (the
  default behavior of System.Text.Json). \\
  \tablerule
  Only add, and do not write defaults & New optional properties with
  \code{[Json\allowbreak{}Ignore(\allowbreak{}Condition =
  When\allowbreak{}Writing\allowbreak{}Default)]} keep existing annotations identical. \\
  \tablerule
  Real format changes only with a major version & Document them, ideally together with the
  EF major upgrade, when users expect migrations anyway. \\
\end{datatable}

\begin{sourcecode}[tests/\allowbreak{}.\allowbreak{}.\allowbreak{}.\allowbreak{}/\allowbreak{}Annotation\allowbreak{}Format\allowbreak{}Tests.\allowbreak{}cs (\allowbreak{}excerpt)]{csharp}
[Fact]
public void Serialized_policy_format_is_stable()
{
    var policy = new RlsPolicy("support_read", PolicyCommand.Select, Permissive: false,
        Roles: ["shop_support"], Using: "true", WithCheck: null);

    // Golden value: changing it is a breaking change for existing snapshots
    Assert.Equal(
        """{"Name":"support_read","Command":"Select","Permissive":false"""
        + ""","Roles":["shop_support"],"Using":"true","WithCheck":null}""",
        policy.Serialize());
}
\end{sourcecode}

Two further tests check tolerant reading in both directions. The snapshot round trip from
chapter~\ref{ch:test-strategy} complements them: the round trip shows that the current code
generator and Fluent API agree, the golden value shows that the next version still agrees
with the snapshots users already have.

% ------------------------------------------------------------
\chapter{Packaging and documentation}\label{ch:packaging-and-documentation}

Shared settings live in \code{Directory.\allowbreak{}Build.\allowbreak{}props}, package
versions centrally in \code{Directory.\allowbreak{}Packages.\allowbreak{}props} (Central
Package Management):

\begin{sourcecode}[Directory.\allowbreak{}Build.\allowbreak{}props (\allowbreak{}shortened)]{xml}
<PropertyGroup>
  <TargetFramework>net10.0</TargetFramework>
  <Nullable>enable</Nullable>
  <ManagePackageVersionsCentrally>true</ManagePackageVersionsCentrally>
  <NoWarn>$(NoWarn);EF1001</NoWarn>
  <Deterministic>true</Deterministic>
</PropertyGroup>
<PropertyGroup>
  <PackageLicenseExpression>MIT</PackageLicenseExpression>
  <PackageReadmeFile>README.md</PackageReadmeFile>
  <VersionPrefix>10.0.0</VersionPrefix>
  <IncludeSymbols>true</IncludeSymbols>
  <SymbolPackageFormat>snupkg</SymbolPackageFormat>
  <PublishRepositoryUrl>true</PublishRepositoryUrl>      <!-- SourceLink -->
  <EmbedUntrackedSources>true</EmbedUntrackedSources>
</PropertyGroup>
\end{sourcecode}

\begin{datatable}{H{0.70}H{1.30}}
  \textbf{Element} & \textbf{Why} \\
  \tablehead
  \code{Generate\allowbreak{}Documentation\allowbreak{}File} & XML docs in the package:
  users see explanations in the IDE. \\
  \tablerule
  \code{Is\allowbreak{}Aot\allowbreak{}Compatible} & Analyzers report AOT and trimming
  problems at build time. \\
  \tablerule
  SourceLink, \code{.\allowbreak{}snupkg} & Users can debug into the plugin's code. Built
  into the SDK since .NET 8. \\
  \tablerule
  README in the package & Shown on nuget.org; minimal example and compatibility. \\
  \tablerule
  CHANGELOG & Mark breaking changes clearly, especially on EF major upgrades. \\
  \tablerule
  \code{Internals\allowbreak{}Visible\allowbreak{}To} & Tests reach internal logic without
  public API. \\
  \tablerule
  Package names & Core without provider, provider packages with a suffix
  (\code{.\allowbreak{}Postgre\allowbreak{}SQL},
  \code{.\allowbreak{}Sql\allowbreak{}Server}). \\
\end{datatable}

Besides the API docs, the documentation above all needs the things users do not see in the
code: which single services are replaced (composability), which event IDs exist, which
limitations apply (pooling) and which security assumptions hold (database roles).

% ------------------------------------------------------------
\chapter{The checklist from A to Z}\label{ch:the-checklist-from-a-to}

For every new plugin, in the order of development:

\section{Architecture}

\begin{checklist}
  \item Decide the level: helpers, model helpers, real plugin or provider?
  \item Separate core and provider packages; types for entity assemblies go into the core.
  \item For every extension point, look for a plugin interface first, replace only then.
  \item Check every mechanism against the simpler alternative with the decision guide
    (chapter~\ref{ch:choosing-the-right-extension-point}).
\end{checklist}

\section{Options}

\begin{checklist}
  \item Extension immutable with copy constructor and virtual Clone.
  \item All model-relevant options in the hash, nothing that varies per context.
  \item Validate for preconditions (provider, combinations).
  \item Maintain LogFragment and PopulateDebugInfo.
  \item Use method in the namespace Microsoft.EntityFrameworkCore, generic and non-generic.
\end{checklist}

\section{Services}

\begin{checklist}
  \item Registration through the services builder (lifetimes from EF's table).
  \item Decorate single services where the public interface suffices
    (chapter~\ref{ch:registering-services}).
  \item Otherwise Replace, and offer the logic publicly and virtually.
  \item ISingletonOptions for singleton services, registered twice.
  \item Classify every piece of state according to chapter~\ref{ch:state-lifetimes-and-thread-safety}: request data never
    in options, hash, model, singleton or translated query.
  \item Reset scoped services with state for pooling (IResettableService).
  \item App dependencies via AddInterceptors, document pooling behavior.
\end{checklist}

\section{Model}

\begin{checklist}
  \item Annotations with a prefix, simple values only, no stored defaults.
  \item Implement the Fluent API non-generically, delegate generically.
  \item Convention builder overloads for ConfigurationSource.
  \item Conventions at ModelFinalizing when they depend on the finished model.
  \item Attributes in the core package, attribute conventions in the provider package.
  \item Validation as IModelFinalizedConvention: throw on errors, warn on questionable
    cases.
\end{checklist}

\section{Type mappings}

\begin{checklist}
  \item Protected parameter constructor and Clone override.
  \item Literal generation for model and provider type.
  \item Public static Default property, shared converter instance.
  \item JsonValueReaderWriter with Instance and ConstructorExpression.
\end{checklist}

\section{Queries}

\begin{checklist}
  \item EF.Functions methods in the core, translation in the provider package.
  \item Translators: null for foreign methods, correct nullability, set type mappings.
  \item IEvaluatableExpressionFilterPlugin for methods without data reference.
  \item Never insert run-time values at translation time (query cache, AOT).
\end{checklist}

\section{Run time}

\begin{checklist}
  \item Implement interceptors synchronously and asynchronously.
  \item Singleton interceptors stateless and thread-safe.
  \item Register singleton interceptors in ApplyServices or pass one shared instance, never a
    new one per context; get per-request data on every call through the context.
  \item Own event IDs in an own range, EventDefinitions cached.
  \item Security-relevant cases as warnings, escalatable via ConfigureWarnings.
  \item Where hooks cannot run, offer an explicit API (e.g. ApplyTenantSession).
\end{checklist}

\section{Migrations \& design time}

\begin{checklist}
  \item Declarative via annotation provider (design time only) and SQL generator.
  \item Escape identifiers, literals via type mappings, validate self-generated SQL.
  \item Custom operations idempotent, extensions in the migrations namespace.
  \item DesignTimeServicesReference with forProvider; code generator removes annotations.
  \item Check compiled models: Default, Instance, runtime annotations.
\end{checklist}

\section{Quality}

\begin{checklist}
  \item Model, SQL, snapshot round trip and integration tests.
  \item Caching test with different options; test for no phantom migration.
  \item Version range per EF major, CI against previews, EF1001 deliberately.
  \item XML docs, SourceLink, symbols, README, CHANGELOG, limitations documented.
  \item Pin serialized annotation formats with golden-value tests.
  \item Test together with popular plugins (naming conventions, decorators) and check the
    coexistence with a second plugin in both orders (chapter~\ref{ch:registering-services}).
\end{checklist}

\section{Advanced (part~\ref{part:VIII})}

\begin{checklist}
  \item Document paths that bypass interceptors (ExecuteUpdate, raw SQL).
  \item Configure value generators with HasValueGeneratorFactory(Type) (compiled models).
  \item Check conventions for complex types and JSON columns.
  \item Document and test the order of custom migration operations.
  \item Consider idempotent scripts and suppressTransaction.
  \item Document operational assumptions (roles, pooling, transaction pooling).
  \item Document context lifetimes: factories, background jobs, open connections.
  \item Cross-tenant access only through separate roles, never through app switches.
  \item Measure hot paths (interceptors), do not guess.
  \item CI with preview build and public API analyzer.
\end{checklist}

\section{Deep extensions (part~\ref{part:IX})}

\begin{checklist}
  \item Custom SQL nodes: VisitChildren, Update, Quote, Print, Equals/GetHashCode.
  \item Handle custom nodes in the SQL generator AND in the nullability processor.
  \item Aggregates: handle predicates, DISTINCT and empty sets deliberately.
  \item Choose runtime annotations deliberately, filter design-time annotations.
  \item Scaffolding: read the plugin's schema objects from the database.
  \item Provider-neutral building blocks into the core, differences into provider packages.
\end{checklist}
